\documentclass{beamer}
\usepackage[utf8]{inputenc}
\usetheme{default}
\title{Protocol Design: Status Update}
\author{Arseny Kurnikov\\
Hasan Islam\\
Harri Hämäläinen}
\begin{document}


\begin{frame}
\titlepage
\end{frame}

\begin{frame}{Current implementation status}
    \begin{itemize}
        \item Implementation in C
        \item Protocol encoding in (BJSON) implemented
        \item Basic protocol mechanisms implemented (subscribe, update,
        keep-alive, ...)
        \item Fragmentation implemented
        \item Transport layer is started
    \end{itemize}
\end{frame}

\begin{frame}{Fragmentation}
    \begin{itemize}
        \item Fragments have the same sequence number
        \item The total number of fragments and the current fragment number
        \item Currently, only the whole message is acknowledged $\Rightarrow$ retransmit
              the whole message if a fragment is lost
    \end{itemize}
\end{frame}

\begin{frame}{Transport layer}
    \begin{itemize}
        \item Packet counters
        \item Sliding window bit rate
        \item NACKs for unreliable connections
        \item Currently one transport layer per subscription $\Rightarrow$ should be one
              transport layer per client
    \end{itemize}
\end{frame}

\begin{frame}{Transport layer 2}
    \begin{itemize}
        \item On-going work
        \item Drift from pure loss-based towards measuring and reporting bit
              rates and losses
        \item The client reports the status of the connection
        \item The server adjusts the sending rate
    \end{itemize}
\end{frame}

\begin{frame}{Metrics}
    \begin{itemize}
        \item Latency
        \item Skew - (TODO) we prioritize low bit-rate flows over high bit-rate
        \item Reliability - remains the same notion of retransmitting only the last update
    \end{itemize}
\end{frame}

\begin{frame}{Law of Jungle}
    \begin{itemize}
        \item Bonald, T.; Feuillet, M.; Proutiere, A., "Is the ``Law of the Jungle''
              Sustainable for the Internet?", INFOCOM 2009.
        \item Theoretical work on teletraffic theory
        \item The assumption of perfect erasure coding or selective retransmission schemes
              that guarantee the usefulness of any received packet
    \end{itemize}

``... there is generally no congestion collapse.  Efficiency remains higher than 90\% ...
as long as maximum source rates are less than link capacity by
one or two orders of magnitude. Moreover, a simple fair drop
policy enforcing fair sharing at flow level is sufficient to guarantee
100\% efficiency in all cases''

\end{frame}

\begin{frame}{Possible design}
    \begin{itemize}
        \item No time to implement and try, just an idea
        \item Current snapshot of sensors at the server, with some buffered data
              for high bit rate sensors
        \item Erasure code them (can prioritize low bit rate sensors)
        \item Send the blocks at some rate
        \item Selectively acknowledge the state
    \end{itemize}
\end{frame}

\end{document}

